% =====================================================================
%  Lisans Bitirme Tezi — Yüz Tanıma Tabanlı Sürücü Güvenlik Sistemi
%  Ege Üniversitesi · Elektrik-Elektronik Mühendisliği · 2025–2026
%
%  Derleme:
%      latexmk -pdf main.tex
%      # veya:  pdflatex main && bibtex main && pdflatex main && pdflatex main
%
%  Türkçe karakter desteği için XeLaTeX gerekmez; pdflatex + babel turkish
%  Gelişme Raporu ile aynı tonu korur.
% =====================================================================
\documentclass[12pt,a4paper]{report}

\usepackage[utf8]{inputenc}
\usepackage[T1]{fontenc}
\usepackage[turkish]{babel}
\usepackage{lmodern}                 % Computer Modern (rapor ile aynı)
\usepackage[a4paper,
            left=3.0cm,right=2.5cm,
            top=2.8cm,bottom=2.5cm]{geometry}
\usepackage{setspace}
\onehalfspacing

\usepackage{graphicx}
\usepackage{caption}
\usepackage{booktabs}
\usepackage{array}
\usepackage{tabularx}
\usepackage{longtable}
\usepackage{microtype}
\usepackage{enumitem}
\usepackage{hyperref}
\usepackage{cleveref}
\usepackage{xcolor}
\usepackage{listings}
\usepackage{tikz}
\usetikzlibrary{shapes.geometric, arrows.meta, positioning, fit, backgrounds}

% --- Renkler (poster ile uyumlu, sober) ----------------------------------
\definecolor{ink}{HTML}{11162A}
\definecolor{accent}{HTML}{8C631A}
\definecolor{rule}{HTML}{2A3047}
\definecolor{mute}{HTML}{5B6075}
\definecolor{signal}{HTML}{9C2A21}

\hypersetup{
    colorlinks=true,
    linkcolor=ink,
    citecolor=accent,
    urlcolor=accent,
    pdftitle={Yüz Tanıma Tabanlı Sürücü Güvenlik ve Acil Durum Bildirim Sistemi},
    pdfauthor={Göktürk Göçen},
    pdfsubject={Lisans Bitirme Tezi, Ege Üniversitesi EE},
}

\captionsetup{font=small, labelfont=bf, labelsep=period}

% --- Listings (kod blokları) ---------------------------------------------
\lstset{
    basicstyle=\ttfamily\footnotesize,
    keywordstyle=\bfseries\color{ink},
    commentstyle=\itshape\color{mute},
    stringstyle=\color{accent},
    showstringspaces=false,
    frame=single,
    rulecolor=\color{rule},
    framesep=4pt,
    breaklines=true,
    breakatwhitespace=true,
    tabsize=2,
    captionpos=b,
}
\lstdefinelanguage{none}{
    keywords={},
    sensitive=false,
}

% --- Türkçe küçük düzeltmeler --------------------------------------------
\addto\captionsturkish{
    \renewcommand{\contentsname}{İçindekiler}
    \renewcommand{\listfigurename}{Şekiller Dizini}
    \renewcommand{\listtablename}{Tablolar Dizini}
    \renewcommand{\bibname}{Kaynakça}
    \renewcommand{\appendixname}{Ek}
    \renewcommand{\chaptername}{Bölüm}
    \renewcommand{\figurename}{Şekil}
    \renewcommand{\tablename}{Tablo}
}

% --- Bölüm başlık biçimlendirme ------------------------------------------
\usepackage{titlesec}
\titleformat{\chapter}[display]
    {\normalfont\Large\bfseries}
    {\chaptertitlename\ \thechapter}{16pt}{\LARGE}
\titlespacing*{\chapter}{0pt}{0pt}{30pt}

% =========================================================================
\begin{document}

% ------------------------------ Kapak ------------------------------------
% ----------------------- Kapak sayfası -----------------------------------
\begin{titlepage}
\begin{center}

{\large\bfseries T.C.}\\[6pt]
{\large Ege Üniversitesi}\\[2pt]
{\large Mühendislik Fakültesi}\\[2pt]
{\large Elektrik--Elektronik Mühendisliği Bölümü}

\vspace{2.2cm}

{\Large\bfseries Lisans Bitirme Projesi}\\[6pt]
{\Large\bfseries Lisans Bitirme Tezi}

\vspace{1.6cm}

{\LARGE\bfseries Yüz Tanıma Tabanlı Sürücü Güvenlik\\[4pt] ve Acil Durum Bildirim Sistemi}\\[10pt]
{\large Driver Safety and Emergency Notification System\\Based on Facial Recognition}

\vspace{1.4cm}

{\bfseries Dönem:} 2025--2026

\vfill

\begin{tabular}{ll}
\bfseries Öğrenci Ad Soyad & Göktürk Göçen \\
\bfseries Öğrenci No       & 05200000XXX  \\
\bfseries E-posta          & qgokturk@gmail.com \\
\bfseries Tez Teslim Tarihi & \today \\
\end{tabular}

\vspace{1.8cm}

\begin{flushleft}
{\bfseries Proje Danışmanı}\\[2pt]
Prof. Dr. Aydoğan Savran\\[14pt]
\hrulefill\\[2pt]
İmza
\end{flushleft}

\end{center}
\end{titlepage}
\cleardoublepage


% ------------------------------ Önsöz / Özet -----------------------------
\frontmatter
\pagenumbering{roman}

\chapter*{Önsöz}
\addcontentsline{toc}{chapter}{Önsöz}

Bu çalışma, Ege Üniversitesi Mühendislik Fakültesi Elektrik--Elektronik
Mühendisliği Bölümü 2025--2026 öğretim yılı Lisans Bitirme Projesi olarak
yürütülmüştür. Proje; taksi gibi araç içi ortamlarda sürücü güvenliğini
artırmaya yönelik, yüz tanıma tabanlı bir kimlik doğrulama ve acil durum
bildirim sistemi tasarımını kapsamaktadır.

Çalışma boyunca değerli rehberliği, sabrı ve dönüş bildirimleriyle proje
danışmanım Sayın Prof. Dr.\ Aydoğan Savran'a teşekkürlerimi sunarım.
Proje, takım arkadaşım Ekin Ağaoğlu ile birlikte iki ayrı tez kapsamında
ortak donanım ve yazılım altyapısı üzerinde geliştirilmiş; ayrıca proje
süresince destek olan aileme ve arkadaşlarıma minnetlerimi sunarım.

\vspace{1cm}
\noindent
\begin{flushright}
Göktürk Göçen \\
\today
\end{flushright}
\cleardoublepage

\chapter*{Özet}
\addcontentsline{toc}{chapter}{Özet}

Bu projede; taksi gibi araç içi ortamlarda sürücü güvenliğini artırmaya
yönelik, gömülü sistem tabanlı bir yüz tanıma çözümü tasarlanmış ve
prototip düzeyinde gerçekleştirilmiştir. Sistem, arka koltuğa yerleştirilen
bir ESP32-CAM modülü üzerinden 2 saniye boyunca 5 kare/saniye hızında
toplam 10 karelik bir görüntü öbeği yakalar ve Wi-Fi üzerinden bulutta
çalıştırılan tanıma servisine iletir. Bulut servisi açık kaynak
InsightFace kütüphanesini (buffalo\_l, ArcFace-R100) kullanarak çoklu
kare için kalite filtresi uygular, geçerli gömme vektörlerinin
centroidini hesaplar ve aranan kişi veri tabanı ile kosinüs benzerliği
ölçer. Sonuç, UART üzerinden STM32 NUCLEO-F767ZI tabanlı olay
yöneticisine iletilir; STM32 LED ve buzzer ile sürücüye geri bildirim
verir ve HM-10 BLE modülü üzerinden eşli Android telefona iletir.
Telefondaki uygulama, BLE bildirimini alarak otomatik biçimde
\textbf{155 Polis İmdat} hattını arar. Tasarım modülerdir; donanım yan~ürünü
olarak elde edilen yerel ön~işleme, bulut tanıma ve BLE tetikleme
katmanları birbirinden bağımsız geliştirilebilir niteliktedir. Çalışma,
düşük maliyetli (donanım eki $<$~500~TL), KVKK uyumlu ve filo araçlarına
genişletilebilir bir referans mimari sunmaktadır.

\noindent\textbf{Anahtar Kelimeler:} yüz tanıma, ArcFace, gömülü sistem,
ESP32-CAM, STM32, BLE, IoT güvenlik, taksi güvenliği, çoklu-kare
agregasyon, pasif canlılık.

\vspace{1cm}

\chapter*{Abstract}
\addcontentsline{toc}{chapter}{Abstract}

This thesis presents the design and prototype implementation of an
embedded face-recognition system aimed at improving driver safety in
in-vehicle environments such as taxis. An ESP32-CAM module mounted toward
the rear seat captures a 10-frame burst at 5\,FPS over 2 seconds upon the
driver's scan command and uploads it over Wi-Fi to a cloud recognition
service. The service, built on top of the open-source InsightFace
library (buffalo\_l, ArcFace-R100) and hosted on AWS EC2, applies a
per-frame quality gate, computes the centroid of valid embeddings, and
performs cosine similarity matching against a wanted-suspect database.
The result is delivered over UART to an STM32 NUCLEO-F767ZI based event
manager that drives the driver-facing LEDs and buzzer and forwards the
event to a paired Android phone over an HM-10 Bluetooth Low Energy link.
The Android application listens for BLE notifications and automatically
dials the national emergency number \textbf{155}. The architecture is
modular: on-device pre-processing, cloud-side recognition and BLE-driven
actuation can be developed independently. The result is a low-cost
(under \euro\,15 in extra hardware), GDPR-aware, fleet-extensible
reference architecture.

\noindent\textbf{Keywords:} face recognition, ArcFace, embedded systems,
ESP32-CAM, STM32, BLE, IoT security, taxi safety, multi-frame aggregation,
passive liveness.

\cleardoublepage


% ------------------------------ İçindekiler ------------------------------
\tableofcontents
\clearpage
\listoffigures
\clearpage
\listoftables
\clearpage

% ------------------------------ Ana metin --------------------------------
\mainmatter
\pagenumbering{arabic}
\setcounter{page}{1}

\chapter{Giriş}
\label{ch:giris}

\section{Konu}

Ticari taksi hizmeti, kentsel ulaşımın temel bileşenlerinden biridir;
buna karşın sürücü ve yolcunun birbirini önceden tanımadığı, kapalı
mekânda ve genellikle gece saatlerinde gerçekleşen bu etkileşim, güvenlik
açısından kırılgan bir noktadır. Sürücüye yönelik saldırı, gasp ve
benzeri olaylar; hem can güvenliği hem de mesleki sürdürülebilirlik
bakımından önemli risk faktörleridir. Mevcut çözümlerin çoğu (panik
butonu, GPS izleme, dijital takograf) ya \emph{olay sonrası} bilgi
sağlamakta ya da sürücünün aktif müdahalesini gerektirmektedir.

Bu çalışmada, sürücünün dikkatini ve sürüş güvenliğini bozmadan,
\textbf{tek bir butona basarak} arka koltuktaki yolcunun aranan kişiler
veri tabanı ile eşleştirilmesini sağlayan, düşük maliyetli ve modüler
bir gömülü sistem prototipi tasarlanmıştır. Sistem; görüntü yakalama,
ön~işleme, bulut tabanlı yüz tanıma, sürücüye görsel/işitsel geri
bildirim ve otomatik 155 çağrısı bileşenlerini uçtan uca entegre eder.

\section{Problem}

Taksi sürücüleri, yolcu kaynaklı tehdit veya saldırı vakalarına karşı
açık bir çalışma ortamına sahiptir. Bu riskleri azaltmak için iki temel
ihtiyaç tanımlanmıştır:

\begin{enumerate}[label=(\roman*),itemsep=2pt]
    \item Yolcu kimliğinin caydırıcılık oluşturacak biçimde
          doğrulanması ve kayıt altına alınması,
    \item Bir olay durumunda hızlı müdahaleyi sağlayacak, sürücünün
          aktif eylemine bağımlı olmayan bir bildirim mekanizmasının
          bulunması.
\end{enumerate}

Var olan çözümler bu iki ihtiyacı tek bir gömülü sistemde, düşük
donanım maliyetiyle ve sürücünün dikkatini dağıtmayacak biçimde bir
araya getirememektedir. Bu eksiklik, mevcut tezin çıkış noktasıdır.

\section{Amaç ve Kapsam}

Tezin amacı, yukarıda tanımlanan iki ihtiyacı karşılayan bir prototip
sistem geliştirmek ve uçtan uca işlevselliğini deneysel olarak
göstermektir. Proje kapsamı aşağıdaki sınırlarla belirlenmiştir:

\begin{itemize}[itemsep=2pt]
    \item Tek aracı kapsayan, sergi ve laboratuvar düzeneğinde test
          edilebilen prototip seviyesinde bir sistem geliştirilecektir.
    \item Kamera ve yüz tanıma; gerçek bir taksi içi montaj senaryosu
          yerine, tavan döşemesi/arka koltuk karşısı ideal koşullar
          baz alınarak değerlendirilecektir.
    \item Acil çağrı yalnızca \textbf{155 Polis İmdat} numarasına
          yönlendirilecek; SMS/e-posta benzeri ek bildirim katmanları
          mimaride yer ayrılmış olsa da bu tezin değerlendirme
          kapsamı dışındadır.
    \item Yolcu görüntüleri sunucu tarafında diske yazılmaz; yalnızca
          işleme süresince RAM'de tutulur. Kalıcı saklama politikası
          KVKK çerçevesinde tartışılır.
\end{itemize}

\section{Hedeflenen Başarı Ölçütleri}
\label{sec:hedefler}

Gelişme raporunda \cite{ekin_gelisme_2026} tanımlanan hedef metriklerden
yola çıkılarak, bitirme tezi kapsamında aşağıdaki başarı ölçütleri
belirlenmiştir:

\begin{itemize}[itemsep=2pt]
    \item \textbf{Uçtan uca tanıma gecikmesi:} ESP32-CAM TARA komutu
          alındıktan, STM32'ye sonucun ulaştığı ana kadar geçen süre
          $\leq$~7~s (10 karelik öbek için, Wi-Fi $\geq$~10\,Mbps yukarı
          hızda).
    \item \textbf{Yüz tanıma doğruluğu (ideal koşullarda):}
          Eşit Hata Oranı (EER) $\leq$~\%5; FAR\,$\leq$\,\%1 noktasındaki
          FRR $\leq$~\%10. Ölçüm \cref{sec:olcum-plani}'ndeki test
          protokolüne göre yapılacaktır.
    \item \textbf{Aydınlatma ve mesafe aralığı:} 100--600~lx aydınlatma ve
          30--120~cm yolcu--kamera mesafesi aralığında kararlı çalışma.
    \item \textbf{Acil çağrı tetikleme süresi:} STM32 MATCH kararından
          telefonun 155'i çevirmeye başlamasına kadar $\leq$~3~s.
    \item \textbf{Ek donanım maliyeti:} Mevcut STM32 kartı dışında
          eklenen bileşenler için $\leq$~500~\textcent\,(TL).
\end{itemize}

\paragraph{Önerilen değişiklik.} Gelişme raporundaki ``uçtan uca
$<$~200~ms'' hedefi, tek karelik bir API sorgusunu varsayıyordu. Çoklu
kare öbek (\emph{burst}) tasarımına geçildiği için bu hedef
\emph{karelerin tamamı için} en fazla 7~s olarak revize edilmiştir
(\cref{sec:degisiklikler}). Yolcu binişi başına yalnızca bir tarama
yapıldığından, kullanıcı deneyimi açısından bu süre uygundur.

\section{Tezin Yapısı}

\Cref{ch:literatur}'de yüz tanıma yöntemleri ve gömülü sistem tabanlı
uygulamalar literatürü özetlenir. \Cref{ch:yontem} sistemin tasarım
ilkelerini ve mimari kararlarını sunar. \Cref{ch:gerceklestirme} ise
gerçekleme aşamasında alınan kararları, kullanılan donanım ve yazılım
bileşenlerini, protokol tanımlarını ve karşılaşılan teknik sorunları
ayrıntılı olarak açıklar. \Cref{ch:sonuclar}, gerçekleştirilen
ölçümleri ve elde edilen sonuçları tartışır. \Cref{ch:sonuc-oneriler}
ise çıkarımları toparlar ve gelecek çalışma önerilerini paylaşır.

\chapter{Literatür Özeti}
\label{ch:literatur}

Bu bölümde, mevcut çalışmanın temellendiği dört eksende literatür
değerlendirilmiştir: (i) yüz tanıma yöntemleri ve performans metrikleri,
(ii) gömülü platformlarda (MCU/SBC) gerçek zamanlı uygulama kısıtları,
(iii) IoT tabanlı acil bildirim sistemleri ve (iv) biyometrik verilerin
yasal/etik çerçevesi.

\section{Yüz Tanıma Yöntemleri ve Performans Metrikleri}

Modern yüz tanıma; tespit (\emph{detection}), hizalama (\emph{alignment})
ve özellik çıkarımı (\emph{embedding}) olmak üzere üç ardışık adımdan
oluşmaktadır. ArcFace mimarisi, son katmandaki softmax fonksiyonuna eklenen
\emph{additive angular margin} terimi sayesinde, sınıflar arası ayrımı
kosinüs uzayında belirginleştirir ve LFW, MegaFace gibi standart veri
kümelerinde insan performansına yakın doğruluk üretir \cite{deng2019arcface}.
ArcFace'in açık kaynak referans uygulaması InsightFace
\cite{guo2021insightface}, RetinaFace \cite{deng2020retinaface} ile
birleştirilmiş hâlde \texttt{buffalo\_l} model paketi olarak
yayınlanmaktadır; bu paket bu çalışmanın bulut katmanında kullanılmıştır.

Yüz tanıma sistemlerinin performansı tipik olarak iki tamamlayıcı oranla
nicelendirilir:

\begin{itemize}
    \item \textbf{Yanlış Kabul Oranı (FAR --- False Accept Rate):} aranan
          listede bulunmayan bir kişinin, aranan biri ile eşleştiği vakaların
          oranı.
    \item \textbf{Yanlış Reddetme Oranı (FRR --- False Reject Rate):}
          gerçek bir aranan kişinin sistem tarafından kaçırıldığı vakaların
          oranı.
\end{itemize}

Bu iki oran, sistemin karar eşiği (kosinüs benzerliği için $\tau$) bir
parametre olarak süpürüldüğünde \emph{karşılıklı ödünleşim} (ROC eğrisi)
şeklinde gözlemlenir. İki oranın eşitlendiği nokta \textbf{Eşit Hata
Oranı (EER --- Equal Error Rate)} olarak adlandırılır ve sistemin tek
sayısal kalite göstergesidir.

Tek bir karelik snapshot ile gerçekleştirilen tanıma, hareket bulanıklığı,
ani aydınlatma değişimi ve geçici tıkanmalardan etkilenir. Birden çok
kare üzerinde gömme vektörlerinin centroid'ini hesaplama yaklaşımı,
literatürde yaygın olarak kullanılan ve tek-shot yaklaşıma kıyasla
gözlemlenebilir doğruluk artışı sağlayan bir tekniktir
\cite{chen2018multiframe}. Bu çalışmada da 10 karelik öbeklenme tercih
edilmiştir.

\section{Sahtelik ve Canlılık (Liveness)}

Yüz tanıma sistemlerinin pratik bir güvenlik açığı, basılı fotoğraf veya
ekran tekrarı (\emph{photo replay}) saldırılarıdır. \emph{Aktif canlılık}
yöntemleri (kullanıcıdan gülümsemesi, gözünü kırpması istenmesi)
caydırıcılığı yüksektir ancak kullanıcı etkileşimi gerektirir; bu da
sürücü için bir dikkat dağıtıcıdır. Bu nedenle, çoklu kare gömme
vektörlerinin kareler arası standart sapması (\emph{cross-frame standard
deviation}) bir \textbf{pasif canlılık} göstergesi olarak kullanılmıştır
\cite{tang2018liveness}. Gerçek bir canlı yüz, mikro-hareketler ve doğal
yüz kasları nedeniyle düşük ama sıfırdan farklı bir standart sapma
üretirken; bir fotoğraf replay senaryosunda standart sapma anlamlı
biçimde azalır. Bu sinyal ek bir model yükü getirmeden, mevcut gömme
vektörleri üzerinden ücretsiz olarak hesaplanmaktadır.

\section{Gömülü Platformlarda Gerçek Zamanlı Uygulama Kısıtları}

Yüz tanıma iş yükünün gömülü tarafa yerleştirilmesi, üç temel kısıt
nedeniyle bu projede tercih edilmemiştir:

\begin{enumerate}
    \item \textbf{Hesaplama yükü:} ArcFace-R100 yaklaşık 24 milyon
          parametreye sahip bir CNN'dir. INT8 nicemleme sonrası bile
          STM32F767 sınıfı bir MCU için (240 MHz, $\leq$~512~KB RAM)
          gecikme bütçesi aşılmaktadır.
    \item \textbf{Bellek kısıtı:} buffalo\_l model paketi, INT8'e
          indirilse dahi 30~MB'ı aşmaktadır; Cortex-M7 MCU'larında
          bu kalıcı bellek sığasının ötesindedir.
    \item \textbf{Geliştirme döngüsü:} Sergi süresi (12 gün) gibi sınırlı
          bir takvimde, hassas nicemleme ve donanım hızlandırma
          çalışmaları öncelikten düşmektedir.
\end{enumerate}

Bu kısıtlar, MCU üzerinde yalnızca \emph{olay yönetimi} (state machine,
LED/buzzer/buton, BLE köprüsü) çalıştırma; tanımayı sunucuya devretme
kararını desteklemektedir. Aynı yaklaşım, Raspberry Pi 4 gibi
tek-kart-bilgisayar (SBC) tabanlı çalışmalarda da
``edge inference + cloud verification'' hibridi olarak
gözlemlenmektedir \cite{kumar2021edgeface}.

\section{IoT Tabanlı Acil Bildirim Sistemleri}

Acil durum tetikleme bileşeni için iki yaygın yaklaşım vardır:
(i) doğrudan gömülü cihazdan GSM modülü ile arama/SMS, (ii) BLE üzerinden
eşli akıllı telefon aracılığı ile arama. İlk yaklaşım, donanım maliyeti
(SIM800/SIM7600 modülleri), SIM kart yönetimi ve operatör onayı gibi ek
yükler getirir. İkinci yaklaşım ise sürücünün zaten yanında bulunan akıllı
telefonu \emph{ağ geçidi} (gateway) olarak kullanır; sistem karmaşıklığını
ve maliyetini düşürür. Bu çalışma ikinci yaklaşımı benimsemiştir.

İki büyük mobil platformun bu konudaki tutumları farklıdır: Android,
\texttt{Intent.ACTION\_CALL("tel:155")} ile arka planda otomatik çağrı
başlatılmasına izin verir \cite{android2024telephony}; iOS ise üçüncü
taraf uygulamalardan tam-otomatik aramaya izin vermez, yalnızca
\texttt{tel://}<numara> URL şemasıyla numara yüklenmiş dialer ekranını
açar \cite{apple_tel_url}. Mevcut tez için, ekibin elinde bulunan iPhone
ile sergiye gidilebilmesi pratik bir tercih olarak iOS'un seçilmesini
sağlamıştır; kullanıcının dialer'da yeşil tuşa tek dokunuşu, hem Apple
sandbox gereği hem de \emph{yanlış pozitif arama koruması} olarak
savunulan bir adım hâline gelmiştir. Android tarafında aynı sistem,
\texttt{Intent.ACTION\_CALL} ile tek-dokunuşsuz çalışacak şekilde
genişletilebilir; bu olasılık \ref{ch:sonuc-oneriler}'te ele alınmıştır.

\section{Biyometrik Verilerin Yasal ve Etik Çerçevesi}

Yüz görüntüsü, KVKK kapsamında \emph{özel nitelikli kişisel veri} olarak
sınıflandırılmaktadır \cite{kvkk2016}. Bu kapsamda toplanan verinin:
(i) açık rıza ile veya kanunlarda öngörülmüş hâllerde işlenmesi,
(ii) işlemenin amacına uygun, asgari ölçüde toplanması,
(iii) öngörülen süre boyunca tutularak sonrasında silinmesi
gerekmektedir. Mevcut tasarımda yolcu görüntüleri yalnızca işleme süresince
RAM'de tutulmakta, sunucunun diskine yazılmamaktadır; sonuç dönüşünden
sonra ilgili oturum yapısı bellekten silinmektedir. Aranan kişiler veri
tabanına yapılacak girişlerin ise yetkili idari/kolluk birimleri
tarafından sağlanması, sistem mimarisinin bir varsayımıdır.

GDPR'nin \emph{purpose limitation} ilkesi de \cite{gdpr2018}, sistemin
yalnızca aranan kişi tespiti amacıyla çalışabileceğini; aynı altyapının
müşteri profillemesi veya benzeri amaçlarla kullanılamayacağını
gerektirir. Tasarımda saklanan tek kalıcı veri, aranan kişilere ait
embedding vektörleridir; ham fotoğraf veri tabanına yazılmaz.

\section{Bu Çalışmanın Konumu}

Yukarıdaki dört eksen değerlendirildiğinde, mevcut çalışmanın özgün
katkıları şu şekilde özetlenebilir:

\begin{itemize}
    \item Açık kaynak InsightFace tabanlı bir tanıma sunucusunun, ESP32-CAM
          ile birlikte uçtan uca taksi senaryosuna uyarlanması;
    \item Çoklu kare centroid agregasyonu ve cross-frame standart sapması
          ile ek model yükü getirmeyen \emph{pasif canlılık} sinyalinin
          birlikte uygulanması;
    \item Sürücü etkileşimini tek butona indiren ve dikkat dağıtmayan bir
          STM32 olay yöneticisi tasarımı;
    \item Akıllı telefonun BLE ağ geçidi olarak kullanılmasıyla, GSM
          modülü gerektirmeyen düşük maliyetli bir acil çağrı zinciri;
    \item Toplam donanım eki maliyetinin 500~TL altında tutulması ve KVKK
          ilkelerine uygun veri yönetimi.
\end{itemize}

\chapter{Yöntem ve Sistem Tasarımı}
\label{ch:yontem}

\section{Sistem Mimarisi}
\label{sec:mimari}

Sistem, fiziksel ve mantıksal olarak üç bölgeden oluşur:

\begin{enumerate}[label=\textbf{\Alph*.},itemsep=2pt]
    \item \textbf{Araç İçi Modül}: ESP32-CAM (kamera + Wi-Fi),
          STM32 NUCLEO-F767ZI (olay yöneticisi), HM-10 BLE köprü modülü,
          harici LED'ler, aktif buzzer, TARA ve PANİK butonları.
    \item \textbf{Bulut Tanıma}: AWS EC2 üzerinde Docker konteynerinde
          çalışan Flask + InsightFace hizmeti, pickle veri tabanı.
    \item \textbf{Acil Çağrı}: Sürücünün BLE üzerinden eşli Android
          telefonu, üzerinde çalışan TaxiGuard uygulaması.
\end{enumerate}

\Cref{fig:sistem-mimarisi} bu üç bölge arasındaki veri akışını gösterir.

\begin{figure}[h!]
\centering
\begin{tikzpicture}[
    node distance=10mm and 10mm,
    box/.style={draw=ink, thick, minimum width=28mm, minimum height=12mm,
                align=center, font=\small\sffamily},
    cloud/.style={draw=ink, thick, minimum width=34mm, minimum height=12mm,
                  align=center, font=\small\sffamily, fill=ink!5},
    arr/.style={-{Latex[length=2mm]}, thick, ink},
    arracc/.style={-{Latex[length=2mm]}, thick, accent},
    arrsig/.style={-{Latex[length=2mm]}, thick, signal},
    lab/.style={font=\scriptsize\ttfamily, ink},
]
\node[box] (esp) {ESP32-CAM\\\scriptsize 10-frame burst, Wi-Fi};
\node[box, below=of esp] (stm) {STM32 F767ZI\\\scriptsize Olay yöneticisi};
\node[box, right=22mm of esp] (wifi) {Wi-Fi 2.4\,GHz};
\node[cloud, right=of wifi] (ec2) {AWS EC2\\\scriptsize Flask + InsightFace};
\node[box, below=of ec2] (phone) {Android\\\scriptsize Intent.ACTION\_CALL};
\node[box, left=of phone] (hm10) {HM-10 BLE};

\draw[arr] (esp.east) -- node[lab, above]{HTTP POST · JPEG} (wifi.west);
\draw[arr] (wifi.east) -- (ec2.west);
\draw[arracc] (ec2.west) to[bend left=30] node[lab, below]{JSON · match} (wifi.east);
\draw[arr] (esp) -- node[lab, right]{UART1\\\scriptsize CAPTURE/RESULT} (stm);
\draw[arr] (stm.east) -- node[lab, above]{UART2} (hm10.west);
\draw[arr] (hm10.east) -- node[lab, above]{BLE notify} (phone.west);
\draw[arrsig] (phone.south) -- ++(0,-7mm) node[lab, right]{GSM · 155};
\end{tikzpicture}
\caption{Sistem üst seviye blok şeması: araç içi modül, bulut tanıma ve acil çağrı katmanları.}
\label{fig:sistem-mimarisi}
\end{figure}

Mimari tercihinde üç temel ilke belirleyici olmuştur:

\begin{itemize}[itemsep=2pt]
    \item \textbf{Tek sorumluluk:} Her bileşen, yalnızca tek bir
          görevden sorumludur (ESP32-CAM: kamera+ağ, STM32: olay yönetimi,
          EC2: tanıma, telefon: çağrı). Bu ayrışma, bir bileşenin
          değişmesinin diğerlerini etkilememesini sağlar.
    \item \textbf{Köprü olarak gömülü:} STM32, kamera ve internet
          erişiminin doğrudan tarafı değildir; ESP32-CAM bu görevi
          devralır. Bu, kamera/Wi-Fi'nin başarısız olması durumunda STM
          tarafındaki olay zincirinin (panik, BLE) etkilenmemesini garanti
          eder.
    \item \textbf{Bulut tarafı içe dönük:} Tanıma sunucusu yalnızca
          tek bir endpoint (\texttt{POST /search}) sunar; tüm karmaşık
          mantık (kalite filtresi, multi-frame agregasyon) burada
          gizlenir, gömülü taraf yalnızca ham JPEG verisi gönderir.
\end{itemize}

\section{Tanıma Yaklaşımı: Alternatiflerin Değerlendirilmesi}
\label{sec:alt-degerlendirme}

Gelişme raporunda \cite{ekin_gelisme_2026} başlatılan ``API destekli mi,
yerel mi?'' tartışması, kapanışında üç alternatifi karşılaştırarak
gerçekleştirme için API destekli yaklaşımı seçmiştir
(\Cref{tab:alternatif-degerlendirme}). Bu seçimin önemli bir nüansı,
``API destekli''nin \emph{kapalı kutu üçüncü taraf bulut servisi}
(AWS Rekognition, Azure Face, Face++) değil; ekibin kendi sunucusuna
deploy ettiği \emph{açık kaynak InsightFace REST API} olarak
yorumlanmasıdır. Bu yorum sayesinde:

\begin{itemize}[itemsep=2pt]
    \item Kullanılan model (\texttt{buffalo\_l}) ve eşik değer
          ($\tau$) ekibin kontrolündedir; tezde FAR/FRR ölçümleri
          dolgun bir savunma için ölçülmüştür.
    \item KVKK uyumu, veri akışını ekibin sahip olduğu sunucuda
          tuttuğu için pratik olarak sağlanmıştır.
    \item Kapatılma riski ya da fiyat değişikliği gibi üçüncü taraf
          bağımlılıkları yoktur.
\end{itemize}

\begin{table}[h!]
\centering
\caption{Yüz tanıma yaklaşımı için alternatif değerlendirme matrisi.}
\label{tab:alternatif-degerlendirme}
\small
\begin{tabularx}{\linewidth}{l X X l}
\toprule
\textbf{Alternatif} & \textbf{Avantaj} & \textbf{Dezavantaj/Risk} & \textbf{Karar}\\
\midrule
MCU üzerinde tanıma (Pico/F767) & Düşük maliyet, internet bağımsız &
Hesap gücü/bellek yetmez; doğruluk düşer & Elendi \\
Yerel ağır model (Pi 4 + CNN) & İnternet bağımsız, gecikme yerel &
Donanım maliyeti, ısı, takvim riski & Ertelendi \\
Kapalı üçüncü-taraf API (AWS Rekognition) & Hızlı entegrasyon &
KVKK belirsizliği, fiyat/kapatma riski, model üzerinde kontrol yok & Elendi \\
\textbf{Kendi InsightFace sunucumuz} & Tam kontrol, KVKK uyumlu, ücretsiz &
Sunucu işletim yükü (EC2 \euro 5/ay) & \textbf{Seçildi} \\
\bottomrule
\end{tabularx}
\end{table}

\section{Çoklu-Kare Burst Akışı}

Tek bir karelik snapshot ile tanıma, geçici tıkanma, hareket bulanıklığı
veya kırpılan yüz açısı durumlarında kırılgan davranır.  Bu projede,
ESP32-CAM tetikleme komutu alır almaz \textbf{5 FPS hızında 2 saniye
boyunca 10 kare} yakalar ve her birini ardışık olarak sunucuya iletir
(\Cref{fig:burst-akisi}). Sunucu tarafı, kareleri tek bir oturum (session)
altında biriktirir; 10.\ kare ulaştığında agregasyonu çalıştırır.

\begin{figure}[h!]
\centering
\begin{tikzpicture}[
    node distance=2mm, font=\scriptsize\ttfamily,
    every node/.style={draw=ink, minimum width=11mm, minimum height=7mm, align=center},
]
\node (n1) {f1\\\scriptsize sid=A,\,i=1/10};
\node (n2) [right=of n1] {f2};
\node (n3) [right=of n2] {f3};
\node (n4) [right=of n3] {f4};
\node (n5) [right=of n4] {f5};
\node (n6) [right=of n5] {f6};
\node (n7) [right=of n6] {f7};
\node (n8) [right=of n7] {f8};
\node (n9) [right=of n8] {f9};
\node (n10) [right=of n9, draw=signal, very thick] {f10\\\scriptsize final};
\end{tikzpicture}
\caption{ESP32-CAM tarafından yakalanan 10 karelik burst. Her kare aynı oturum kimliğini taşır;
yalnızca son kareye sunucu \emph{karar} ile yanıt verir, önceki karelerin yanıtı bir devam (continue) bildirimidir.}
\label{fig:burst-akisi}
\end{figure}

Bu yaklaşımın iki kazancı vardır:

\begin{enumerate}[label=(\arabic*),itemsep=2pt]
    \item \textbf{Robustluk:} 10 kareden bir kısmı kalite filtresine
          takılsa bile, en az 5 geçerli kare bulunduğunda doğruluk
          tek-kare durumuna kıyasla belirgin biçimde yükselir
          (literatür gözlemi: \cite{chen2018multiframe}).
    \item \textbf{Pasif canlılık:} 10 karelik gömme vektörlerinin
          karelere göre standart sapması, \emph{cross-frame embedding
          variance} olarak hesaplanır ve fotoğraf replay saldırılarını
          ek bir model yükü olmadan tespit etmek için kullanılır.
\end{enumerate}

\section{Sunucu Tarafı Kalite Filtresi ve Agregasyon}

Sunucu, her kare için aşağıdaki kalite kapısını uygular:

\begin{itemize}[itemsep=1pt]
    \item Yüz tespit skoru (\texttt{det\_score}) $\geq$ 0.7,
    \item Yüz bölgesi alanı (piksel) $\geq$ 80$\times$80,
    \item Yüz dönüklüğü (yaw) açısı $|$yaw$|$ $\leq$ 30$^{\circ}$,
    \item Laplace varyansı tabanlı bulanıklık skoru $\geq$ 10.
\end{itemize}

Bu kapıyı geçen karelerin embedding vektörleri biriktirilir. Son karede
en az \texttt{MIN\_QUALITY\_FRAMES} (varsayılan 5) geçerli kare yoksa,
sonuç \texttt{insufficient\_quality\_frames} olarak işaretlenir; ESP'ye
NOMATCH yerine kalite-yetersiz uyarısı döner. Yeterli kare varsa:

\begin{equation}
\vec{c} = \frac{1}{N}\sum_{i=1}^{N} \vec{e}_i, \quad
\vec{e}_i = \mathrm{ArcFace}(\mathrm{crop}(I_i)),
\end{equation}

şeklinde centroid hesaplanır ve veri tabanındaki her kayıt $\vec{d}_k$
ile kosinüs benzerliği ölçülür:

\begin{equation}
\mathrm{sim}_k = \frac{\vec{c} \cdot \vec{d}_k}{\|\vec{c}\|\, \|\vec{d}_k\|}.
\end{equation}

En yüksek benzerlik $\mathrm{sim}^\star = \max_k \mathrm{sim}_k$, eşik
$\tau$ (varsayılan 0.40) ile karşılaştırılır:
$\mathrm{sim}^\star \geq \tau \Rightarrow \text{MATCH}$. Pasif canlılık
skoru ise:

\begin{equation}
\mathrm{liveness} = \frac{1}{D}\sum_{j=1}^{D} \sigma_j, \quad
\sigma_j = \mathrm{std}\bigl(\{e_{i,j}\}_{i=1}^{N}\bigr),
\end{equation}

ile, embedding boyutlarına göre standart sapmaların ortalaması olarak
hesaplanır. Gerçek bir canlı yüzde bu skor sıfıra yakın ama anlamlı
biçimde sıfırın üstündedir; fotoğraf replay durumunda anlamlı biçimde
azalır.

\section{STM32 Olay Yöneticisi}
\label{sec:stm-olay-yonetimi}

STM32 NUCLEO-F767ZI üzerinde çalışan firmware, bir sonlu durum makinesi
(\emph{finite state machine}, FSM) etrafında düzenlenmiştir.
Modül HAL bağımlılığından arındırılmış, taşınabilir C99 olarak yazılmış
ve donanımla iletişimi bir \emph{callback vtable} üzerinden
gerçekleştirmektedir (\Cref{fig:fsm}). Bu sayede aynı kaynak dosyası
hem CubeIDE projesinde derlenebilmekte, hem de PC üzerinde sahte saat
ve sahte callback'lerle birim testten geçirilebilmektedir.

\begin{figure}[h!]
\centering
\begin{tikzpicture}[
    node distance=18mm and 18mm,
    state/.style={draw=ink, thick, rounded corners=2pt,
                  minimum width=20mm, minimum height=9mm,
                  align=center, font=\small\sffamily},
    accent/.style={draw=signal, very thick},
    arr/.style={-{Latex[length=1.6mm]}, thick, ink, font=\scriptsize\ttfamily},
]
\node[state] (idle) {IDLE};
\node[state, right=of idle] (scan) {SCANNING};
\node[state, above right=8mm and 14mm of scan] (match) {MATCH};
\node[state, below right=8mm and 14mm of scan] (nom) {NOMATCH};
\node[state, accent, below=20mm of scan] (panic) {PANIC};
\node[state, right=of nom] (neterr) {NETERR};

\draw[arr] (idle) -- node[above]{TARA} (scan);
\draw[arr] (scan) -- node[sloped, above]{MATCH} (match);
\draw[arr] (scan) -- node[sloped, below]{NOMATCH} (nom);
\draw[arr] (scan) -- node[sloped, above]{timeout / ERR} (neterr);
\draw[arr, dashed] (match.west) to[bend left=30] node[above]{hold 5\,s} (idle.north);
\draw[arr, dashed] (nom.west) to[bend right=30] node[below]{hold 2\,s} (idle.south);
\draw[arr, dashed] (neterr.north) to[bend left=15] node[right]{hold 3\,s} (idle.east);
\draw[arr, signal] (idle.south) to[bend right=25] node[left]{PANIC} (panic.west);
\draw[arr, signal] (scan.south) to[bend left=15] node[right]{PANIC} (panic.east);
\draw[arr, dashed] (panic.west) to[bend left=45] node[left]{hold 10\,s} (idle.south);
\end{tikzpicture}
\caption{STM32 olay yöneticisi durum geçiş diyagramı. PANIC olayı her durumda erişilebilir.}
\label{fig:fsm}
\end{figure}

Durum çıktıları aşağıdaki gibi tanımlanmıştır:

\begin{itemize}[itemsep=1pt]
    \item \textbf{IDLE:} yeşil LED sabit, 5\,s'de bir \texttt{HB\textbackslash n} kalp atışı,
    \item \textbf{SCANNING:} sarı LED sabit, UART1'e \texttt{CAPTURE\textbackslash n},
          UART2'ye \texttt{SCANNING\textbackslash n},
    \item \textbf{MATCH:} kırmızı LED + buzzer, HM-10'a \texttt{MATCH:<name>;<sim>\textbackslash n},
    \item \textbf{NOMATCH:} yeşil LED, HM-10'a \texttt{NOMATCH\textbackslash n},
    \item \textbf{PANIC:} kırmızı LED + buzzer, HM-10'a \texttt{PANIC\textbackslash n},
    \item \textbf{NETERR:} sarı LED yanıp sönüyor, HM-10'a \texttt{NETERR\textbackslash n}.
\end{itemize}

\section{Acil Çağrı Zinciri: BLE ve Android}

HM-10 modülü \emph{transparent UART$\leftrightarrow$BLE} köprüsü olarak
çalışır: STM32'nin UART2'ye yazdığı her bayt, modülün GATT karakteristiği
(\texttt{FFE0} servisi, \texttt{FFE1} karakteristiği) üzerinden eşli
telefona \emph{notification} olarak iletilir. Telefon tarafındaki TaxiGuard
uygulaması bir \emph{foreground service} olarak BLE central rolünde
durmakta, gelen ASCII satırları parse etmekte ve \texttt{MATCH:..} veya
\texttt{PANIC} olayı geldiğinde \texttt{Intent.ACTION\_CALL("tel:155")}
ile otomatik olarak Polis İmdat hattını aramaktadır.

\Cref{tab:mesaj-cercevesi} sistemde tanımlı mesaj çerçevelerini özetler.

\begin{table}[h!]
\centering
\caption{STM32 $\rightarrow$ HM-10 $\rightarrow$ Android mesaj çerçeveleri.}
\label{tab:mesaj-cercevesi}
\small
\begin{tabular}{l l l}
\toprule
\textbf{Mesaj} & \textbf{Tetik} & \textbf{Telefon eylemi}\\
\midrule
\texttt{SCANNING\textbackslash n} & TARA başlangıcı     & Banner: ``Tarama başladı''\\
\texttt{MATCH:<name>;<sim>\textbackslash n} & Pozitif eşleşme & \texttt{Intent.ACTION\_CALL("tel:155")}\\
\texttt{NOMATCH\textbackslash n}  & Burst negatif       & Banner: ``Eşleşme yok''\\
\texttt{PANIC\textbackslash n}    & PANİK butonu        & \texttt{Intent.ACTION\_CALL("tel:155")}\\
\texttt{NETERR\textbackslash n}   & Scan timeout/ERR    & Banner: ``Sunucuya erişilemedi''\\
\texttt{HB\textbackslash n}       & 5\,s aralıkla IDLE'da & Bağlantı sağlamlık göstergesi\\
\bottomrule
\end{tabular}
\end{table}

iOS yerine Android'in seçilme nedeni, iOS sandbox kısıtlamasının
\texttt{ACTION\_CALL} dengini desteklememesidir; iOS yalnızca
dialer'ı önceden doldurulmuş hâlde açabilen
\texttt{tel:} URL şemasına izin vermekte, otomatik bağlantı için
kullanıcı dokunması gerekmektedir.

\section{Ağ Yokken Davranış (Safe Mode)}
\label{sec:safe-mode}

Wi-Fi'nin yokluğunda veya sunucunun yanıt vermediği durumda sistem
\textbf{güvenli moda} geçer:

\begin{enumerate}[label=(\arabic*),itemsep=1pt]
    \item ESP32-CAM, HTTP POST sırasında 10~s zaman aşımına ulaşırsa
          UART1'e \texttt{ERR:no\_wifi} satırı gönderir.
    \item STM32 olay yöneticisi 15~s'lik scan timeout'una takılırsa
          (veya \texttt{ERR:..} satırı alırsa) NETERR durumuna geçer.
    \item NETERR durumunda sarı LED yanıp söner, buzzer susar ve HM-10'a
          \texttt{NETERR\textbackslash n} satırı gönderilir; telefon
          ``Sunucuya erişilemedi'' bildirimiyle sürücüyü uyarır.
    \item Bağlantı geri geldiğinde sistem normal akışa döner; ek bir
          yeniden bağlanma adımı gerektirmez.
\end{enumerate}

Sistem, ağ yokken sessizce başarısız olmaz; sürücüye görsel ve metinsel
geri bildirim ile durumu açıkça bildirir. Bu, hedef başarı ölçütleri
arasında tanımlanan ``ağ olmadığında güvenli mod'' gerekliliğini
karşılar.

\section{Öneri Raporuna ve Gelişme Raporuna Göre Değişiklikler}
\label{sec:degisiklikler}

Bu tezde, gelişme raporundaki tasarımdan iki önemli sapma vardır:

\begin{itemize}[itemsep=2pt]
    \item \textbf{Kamera platformu:} Plan A'da STM32'nin DCMI çevre
          birimine bağlı stand-alone bir OV5640 modülü öngörülmüştü.
          Sergi takvimi içinde DVP/parallel arabirimli OV5640 modülü
          tedarik edilemediği için \textbf{Plan B} olarak ESP32-CAM
          (AI-Thinker, OV3660 sensörü) kullanılmıştır. Bu değişikliğin
          \emph{tasarım üzerindeki etkisi:} kamera + Wi-Fi tek board'a
          taşınmış, STM32 yalnızca olay yöneticisi rolünde kalmıştır.
    \item \textbf{Gecikme hedefi:} Gelişme raporundaki ``$<$~200~ms''
          tek-kare ile çalışan bir API sorgusunu varsayıyordu. Çoklu kare
          burst tasarımına geçildiği için bu hedef, 10 kare için
          \emph{en fazla 7~s} olarak revize edilmiştir. Bu süre,
          yolcu binişi başına bir kez gerçekleşen bir işlem olduğu için
          kullanıcı deneyimi açısından uygundur.
\end{itemize}

\chapter{Gerçekleştirme}
\label{ch:gerceklestirme}

Bu bölümde, \cref{ch:yontem}'de tanımlanan tasarımın somut donanım ve
yazılım bileşenlerine indirgenişi anlatılmaktadır. Her alt sistem için
\emph{ne kullanıldı, neden kullanıldı, nasıl bağlandı, hangi
parametrelerle çalıştığı} sırasıyla verilmiştir.

\section{Donanım Listesi}

\Cref{tab:bom} prototipte kullanılan tüm donanım bileşenlerini ve
yaklaşık birim maliyetlerini özetler. Toplam ek donanım maliyeti,
hedeflenen 500~TL bütçesi içinde kalmıştır.

\begin{table}[h!]
\centering
\caption{Bileşen listesi (BOM) ve yaklaşık birim maliyetler.}
\label{tab:bom}
\small
\begin{tabular}{l l c l}
\toprule
\textbf{Parça} & \textbf{Model / Tip} & \textbf{Adet} & \textbf{Tahmini Fiyat}\\
\midrule
Olay yöneticisi MCU      & STM32 NUCLEO-F767ZI      & 1 & (Mevcut, bütçeye eklenmedi)\\
Kamera + Wi-Fi modülü    & ESP32-CAM AI-Thinker (OV3660) & 1 & 120 TL\\
BLE köprü modülü         & HM-10 (CC2541)           & 1 & 80 TL\\
LED                      & Yeşil 5\,mm, sarı, kırmızı + 220\,$\Omega$ direnç & 3+3 & 10 TL\\
Buzzer                   & Aktif buzzer 5\,V        & 1 & 15 TL\\
Buton                    & Push button taktil 6\,mm & 2 & 10 TL\\
Breadboard + jumper      & --                       & -- & 100 TL\\
USB-TTL dönüştürücü      & CP2102 / FT232           & 1 & 50 TL\\
Yedek 40-pin male header & --                       & 2 & 10 TL\\
\midrule
\textbf{Toplam ek donanım} & & & \textbf{$\approx$ 395 TL}\\
\bottomrule
\end{tabular}
\end{table}

Buluttaki çalışma maliyeti için AWS EC2 \texttt{m7i-flex.large}
örneği eu-central-1 (Frankfurt) bölgesinde çalıştırılmaktadır;
AWS Başlangıç Kredisi ile sergi ve test dönemi için ek maliyet
oluşturmamıştır.

\section{Sunucu Tarafı: Flask + InsightFace}
\label{sec:sunucu-impl}

\subsection{Yığın ve Bağımlılıklar}

Sunucu tarafı Python 3.11 ile yazılmış, Flask web framework'ü ve
InsightFace 0.7 paketi üzerine kurulmuştur. ONNX Runtime CPU sürümü
kullanılmaktadır; bu sayede üretim ortamı (Linux/EC2) ve geliştirme
ortamı (macOS) arasında deterministik davranış sağlanmaktadır.
Konteynerleştirme Docker ile yapılmış; Gunicorn $-w\,1$ yapılandırmasıyla
çalıştırılmaktadır. Tek worker zorunluluğu, oturum (session) verisinin
süreç içi (in-memory) saklanmasından kaynaklanır: birden fazla worker,
aynı oturumun karelerini farklı süreçlere dağıtarak agregasyonu kırardı.

\subsection{Endpoint Tasarımı}

Tek bir kritik endpoint mevcuttur: \texttt{POST /search}. İstek
gövdesi JPEG bayt dizisi; başlıkları:

\begin{itemize}[itemsep=1pt]
    \item \texttt{X-Session-Id}: Oturum kimliği (UUID), tüm burst boyunca aynı.
    \item \texttt{X-Frame-Index}: 1\textendash N arası kare sırası.
    \item \texttt{X-Frame-Total}: Burst içindeki toplam kare sayısı (\=10).
    \item \texttt{X-Shared-Secret}: (Opsiyonel) İlk faz için kapalı bir
          anahtar; üretimde set edildiğinde \texttt{X-Shared-Secret}'i
          taşımayan istekler 401 ile reddedilir.
\end{itemize}

Ara karelere yanıt: \texttt{\{"status":"continue", ...\}}; son kareye
yanıt:

\begin{lstlisting}[language=none, caption={Son kareye sunucu yanıt örneği.}, label=lst:final-response]
{
  "match": true,
  "name": "Ali_Yilmaz",
  "similarity": 0.9437,
  "frames_used": 10,
  "frames_total": 10,
  "liveness_score": 0.014823
}
\end{lstlisting}

\subsection{Operasyonel Sertleştirmeler}

Sergi öncesi yapılan üç hafif sertleştirme:

\begin{enumerate}[itemsep=2pt]
    \item \textbf{Opsiyonel paylaşılan anahtar:} \texttt{SHARED\_SECRET}
          ortam değişkeni atandığında \texttt{/search} endpoint'i
          \texttt{X-Shared-Secret} başlığını arar; \texttt{/health} ve
          \texttt{/metrics} izleme için açık tutulur.
    \item \textbf{Yapılandırılmış JSON günlüğü:} Her burst için bir
          satırlık JSON kayıt (\texttt{event:"burst\_complete"}, oturum
          kimliği, match kararı, similarity, liveness, gecikme) standart
          çıktıya yazılır. Bu format, sergi sonrası analiz ve tez Bölüm 5
          tablolarının üretimi için doğrudan kullanılabilir.
    \item \textbf{\texttt{/metrics} endpoint'i:} Açık oturum sayısı,
          tamamlanan burst sayısı, MATCH/NOMATCH dağılımı, son 100
          burst'ün gecikme p50/p95/maks değerleri, uptime ve kimlik
          doğrulama durumu tek bir JSON cevabıyla raporlanır.
\end{enumerate}

\section{ESP32-CAM: Kamera ve Ağ Köprüsü}

\subsection{Donanım Yapılandırması}

AI-Thinker ESP32-CAM modülü, ESP32-WROOM-32 SoC üzerine entegre OV3660
kamera modülü ve antenden oluşur. Wi-Fi STA modunda sürücünün telefon
hotspot'una bağlanır; flash etmek için 3.3~V FTDI bağlantısı kullanılır
(IO0 GND'ye çekilerek bootloader moduna geçilir).

\subsection{Yazılım Akışı}

Firmware PlatformIO + Arduino-ESP32 çekirdeği üzerinde C++ olarak
yazılmıştır. Ana döngü iki görevi yürütür:

\begin{enumerate}[label=(\arabic*),itemsep=1pt]
    \item STM32'den UART üzerinden gelen \texttt{CAPTURE\textbackslash n}
          komutunu dinlemek,
    \item Komut geldiğinde 5~FPS hızında 10 JPEG kare yakalamak, her
          birini yeni bir HTTP POST isteğiyle EC2'ye iletmek ve son
          karenin JSON cevabını parse etmek.
\end{enumerate}

Burst boyunca her POST aynı oturum kimliğini (\texttt{X-Session-Id})
taşır. Son kare üzerindeki yanıt JSON olarak parse edildikten sonra:

\begin{itemize}[itemsep=1pt]
    \item Eşleşme varsa: STM32'ye \texttt{1;<name>;<sim>\textbackslash n} CSV satırı,
    \item Eşleşme yoksa: \texttt{0;;<sim>\textbackslash n},
    \item Hata durumunda: \texttt{ERR:<tag>\textbackslash n} (örnek tag:
          \texttt{no\_wifi}, \texttt{http\_500}, \texttt{json\_parse}).
\end{itemize}

\subsection{Yapılandırma}

Wi-Fi bilgileri, sunucu URL'si ve TLS bayrağı \texttt{include/config.h}
dosyasından okunur. Faz 1'de (HTTP) ve Faz 2'de (HTTPS + Let's Encrypt)
arasındaki tek farklılık bu dosyadır; sertifika güncellemesi ve TLS
bayrağı ayarı sonrasında yeni image flash edilir.

\section{STM32 Olay Yöneticisi: Gerçekleme}

\subsection{Donanım Konfigürasyonu}

CubeMX projesi NUCLEO-F767ZI üzerinde aşağıdaki çevre birimleri ile
yapılandırılmıştır:

\begin{itemize}[itemsep=1pt]
    \item UART1, 115200~bps 8N1, PA9/PA10: ESP32-CAM ile haberleşme,
    \item UART2, 9600~bps 8N1, PD5/PD6: HM-10 ile haberleşme,
    \item GPIO: PB0 (yeşil), PD12 (sarı), PB14 (kırmızı), PD13 (buzzer),
    \item EXTI: PC13 (TARA butonu, kart üstü mavi USER butonu), PA0 (harici PANİK butonu).
\end{itemize}

\textbf{Önemli not:} F767ZI'de Ethernet çevre birimi varsayılan olarak
aktif gelir ve DCMI pinleri ile çakışır. Bu projede DCMI kullanılmadığı
için ETH disabled bırakılmıştır; eski Plan A taslağındaki kamera doğrudan
DCMI'ye bağlanacak senaryosu Plan B'ye geçişle birlikte düşmüştür.

\subsection{Yazılım Mimarisi}

STM32 firmware'i iki katmanlı geliştirilmiştir.

\paragraph{(i) Tasarım doğrulama katmanı.}
Durum makinesi mantığı, taşınabilir C99 olarak yazılmış ayrı bir
\textbf{referans modül} (\texttt{state\_machine.c/.h}) hâlinde ifade
edilmiş ve PC üzerinde birim testleriyle doğrulanmıştır. Modül, HAL'a
doğrudan bağımlı değildir; tüm donanım etkileşimleri bir callback
vtable üzerinden gerçekleşir:

\begin{lstlisting}[language=C, caption={State machine referans modülünün callback vtable'ı.}, label=lst:cb]
typedef struct {
    void     (*set_led)(sm_led_t led, bool on);
    void     (*set_buzzer)(bool on);
    void     (*send_to_esp)(const char *line);
    void     (*send_to_hm10)(const char *line);
    uint32_t (*now_ms)(void);
} sm_callbacks_t;
\end{lstlisting}

Bu soyutlama sayesinde modülün davranışı donanım olmadan
gcc/clang ile derlenen bir host-side test koşumunda deterministik olarak
doğrulanmıştır. \Cref{tab:host-tests} kapsanan 12 senaryoyu özetler.

\paragraph{(ii) Embedded gerçekleme.}
Tasarım doğrulandıktan sonra, \texttt{main.c} dosyasında aynı durum
makinesi mantığı CubeMX'in yeniden ürettiği \emph{USER CODE} bölgelerine
\emph{inline} edilmiştir. Bu pratik tercih iki nedenden ötürü
yapılmıştır: (i) CubeMX yenileme süreci ile çakışma riskini azaltmak,
(ii) embedded tarafta UART RX kesme rutininin ve EXTI callback'lerinin
state machine eventleri ile aynı dosyada yaşamasının okunabilirliği
artırması. Embedded ve referans uygulamalar arasındaki davranış birebir
aynıdır; tek farklılık hold timer sabitleri (NOMATCH 1\,s vs.\ 2\,s
gibi), proje yürüyüşü sırasında pratiğe göre kalibre edilmiştir.

Bu iki katmanlı yaklaşım, donanım olmadan tasarım hatalarını yakalama
ve donanım üzerinde flash--flash döngülerini en aza indirme açısından
proje takvimini önemli ölçüde rahatlatmıştır.

\begin{table}[h!]
\centering
\caption{STM32 durum makinesi için host-side birim test kapsamı.}
\label{tab:host-tests}
\small
\begin{tabular}{l l}
\toprule
\textbf{Test senaryosu} & \textbf{Doğrulanan davranış}\\
\midrule
initial state              & IDLE + yeşil LED sabit\\
TARA press                 & SCANNING geçişi + CAPTURE komutu\\
ESP MATCH                  & MATCH durumu + kırmızı LED + buzzer + HM-10 frame\\
ESP NOMATCH                & NOMATCH yolu, yeşil LED, doğru frame\\
scan timeout               & 15~s sonrası NETERR\\
PANIC during scan          & PANIC durumunun her durumdan erişilebilir olması\\
hold expiry                & MATCH/NOMATCH/NETERR sonunda IDLE'a dönüş\\
heartbeat                  & IDLE'da 5~s aralıkla HB satırı\\
TARA during hold           & Hold penceresinin kesilebilmesi\\
late ESP result            & Timeout sonrası gelen yanıtın yok sayılması\\
NETERR blink               & NETERR'de sarı LED'in periyodik geçişi\\
oversized name             & 32 karakterden uzun isimde taşma yok\\
\bottomrule
\end{tabular}
\end{table}

\subsection{Embedded Uygulama --- main.c}

\texttt{main.c} dosyasında durum makinesi mantığı USER CODE bölgelerine
inline yazılmıştır. \Cref{lst:adapter} temel parçaları gösterir.

\begin{lstlisting}[language=C, caption={CubeIDE main.c içerisindeki durum makinesi gerçeklemesi (özet).}, label=lst:adapter]
typedef enum {
    STATE_IDLE = 0, STATE_SCANNING, STATE_MATCH,
    STATE_NOMATCH,  STATE_PANIC,    STATE_NETERR,
} app_state_t;

static volatile app_state_t app_state = STATE_IDLE;
static uint32_t state_entered_ms = 0;

static void apply_state_outputs(app_state_t s) {
    /* tüm LED + buzzer kapat, sonra duruma göre aç */
    HAL_GPIO_WritePin(LED_GREEN_GPIO_Port,  LED_GREEN_Pin,  GPIO_PIN_RESET);
    /* ... */
    switch (s) {
        case STATE_IDLE: HAL_GPIO_WritePin(LED_GREEN_GPIO_Port, LED_GREEN_Pin, GPIO_PIN_SET); break;
        case STATE_MATCH: case STATE_PANIC:
            HAL_GPIO_WritePin(LED_RED_GPIO_Port, LED_RED_Pin, GPIO_PIN_SET);
            HAL_GPIO_WritePin(BUZZER_GPIO_Port,  BUZZER_Pin,  GPIO_PIN_SET);
            break;
        /* ... */
    }
}

void HAL_GPIO_EXTI_Callback(uint16_t pin) {
    if (pin == BTN_SCAN_Pin)  flag_tara  = true;
    if (pin == BTN_PANIC_Pin) flag_panic = true;
}
\end{lstlisting}

UART1 RX yolu (ESP'den gelen \texttt{RESULT:..}/\texttt{ERR:..}/\texttt{HB}
satırları), bir bayt-bayt kesme rutini ile \texttt{\textbackslash n}'a kadar bir
satır tamponuna birikir; ana döngüde tamamlanan satırlar
\texttt{process\_cam\_line()} ile parse edilerek MATCH/NOMATCH/NETERR
geçişlerini tetikler.

\section{Android Uygulaması: TaxiGuard v2}

\subsection{Bileşenler}

Uygulama Kotlin + Jetpack Compose ile yazılmış, hedef SDK 34, minimum
SDK 26'dır. Üç ana modülden oluşur:

\begin{itemize}[itemsep=2pt]
    \item \textbf{\texttt{Hm10Service}}: Bir foreground servis olarak
          BLE central rolünde çalışır. HM-10 modülünü \texttt{0xFFE0}
          servis UUID'sine göre tarar, bulunduğunda GATT bağlantısı kurar
          ve \texttt{0xFFE1} karakteristiğinin notifikasyonlarına abone
          olur. MTU parçalanmasına karşı dahili bir \emph{line buffer}
          her \texttt{\textbackslash n}'a kadar baytları biriktirir.
    \item \textbf{\texttt{FrameParser}}: Gelen ASCII satırları parse
          ederek bir \texttt{MatchEvent} sealed hierarchy'sine dönüştürür
          (\texttt{Scanning, Match(name, sim), NoMatch, Panic, NetError,
          Heartbeat, Unknown}).
    \item \textbf{\texttt{EmergencyDialer}}: MATCH veya PANIC event'i
          alındığında \texttt{Intent.ACTION\_CALL("tel:155")}'i çağırır;
          \texttt{CALL\_PHONE} izni yoksa \texttt{ACTION\_DIAL} dönüşüne
          fallback yapar (dialer önceden doldurulmuş şekilde açılır).
\end{itemize}

UI tek ekrandır (\texttt{HomeScreen}); bağlantı bandı, son MATCH kartı,
ekran içi olay günlüğü ve servis Başlat/Durdur kontrolleri içerir.
Ayrıca, donanım yokken bile parse + dial akışını tek telefonda doğrulamak
için \emph{``DEV: Simulate MATCH''} butonu yer alır. Bu buton, servise
\texttt{ACTION\_SIMULATE\_MATCH} intent'i göndererek sahte bir
\texttt{MATCH:Test\_Subject;0.95} satırını parse pipeline'ına besler;
böylece HM-10 elimizde olmadan da uygulamanın MATCH$\rightarrow$155
zinciri doğrulanabilmektedir.

\subsection{İzinler ve Sandbox Davranışı}

\Cref{tab:android-izinler} uygulamanın manifest'inde tanımlı izinleri ve
gerekçelerini özetler.

\begin{table}[h!]
\centering
\caption{Android izinleri ve gerekçeleri.}
\label{tab:android-izinler}
\small
\begin{tabular}{l l}
\toprule
\textbf{İzin} & \textbf{Gerekçe}\\
\midrule
\texttt{BLUETOOTH\_SCAN}                          & API 31+ BLE cihaz keşfi\\
\texttt{BLUETOOTH\_CONNECT}                        & GATT bağlantı kurma\\
\texttt{CALL\_PHONE}                               & \texttt{Intent.ACTION\_CALL("tel:155")}\\
\texttt{FOREGROUND\_SERVICE\_CONNECTED\_DEVICE}    & Servis kategorisi (BLE)\\
\texttt{POST\_NOTIFICATIONS}                       & API 33+ foreground bildirimi\\
\texttt{ACCESS\_FINE\_LOCATION}                    & API 26--30 (legacy BLE keşif)\\
\bottomrule
\end{tabular}
\end{table}

\section{Ölçüm Harness'i: \texttt{far\_frr.py}}
\label{sec:harness}

Tez \cref{ch:sonuclar} ölçümlerinin tekrarlanabilir biçimde
yapılabilmesi için \texttt{eval/far\_frr.py} adlı bir Python aracı
geliştirilmiştir. Araç bir veri seti klasörünü (\texttt{dataset/<isim>/<foto>.jpg})
tarar, her fotoğrafı 10 karelik bir burst olarak \texttt{/search}
endpoint'ine gönderir ve dönen similarity skorlarını CSV'ye yazar.
Ardından eşik $\tau \in [0, 1]$ üzerinde 0.01 adımlarla bir süpürme
gerçekleştirerek FAR, FRR ve TA/FA/TR/FR sayılarını üretir; Equal Error
Rate (EER) eşiğini ve istenen maksimum FAR altındaki en iyi operating
point'i raporlar. Matplotlib mevcutsa bir FAR--FRR grafiği de
oluşturur.

\Cref{tab:harness-cikti} aracın ürettiği çıktı dosyalarını özetler.

\begin{table}[h!]
\centering
\caption{\texttt{far\_frr.py} aracının ürettiği çıktı dosyaları.}
\label{tab:harness-cikti}
\small
\begin{tabular}{l l}
\toprule
\textbf{Dosya} & \textbf{İçerik}\\
\midrule
\texttt{<run>.csv}            & Her probe için bir satır; similarity, matched name, kategori\\
\texttt{<run>\_sweep.csv}     & Eşik--FAR/FRR tablosu (101 satır)\\
\texttt{<run>\_summary.json}  & EER eşiği, operating point, sayımlar\\
\texttt{<run>\_roc.png}       & FAR--FRR eğrisi (matplotlib varsa)\\
\bottomrule
\end{tabular}
\end{table}

Aracın ürettiği CSV ve JSON dosyaları, tezin tablo ve grafiklerine
doğrudan kaynak olur; \texttt{run\_yyyy\_mm\_dd} isimlendirmesi sayesinde
hangi koşulda yapılmış ölçüm olduğu izlenebilir.

\chapter{Ölçüm Düzeni ve Sonuçlar}
\label{ch:sonuclar}

\section{Ölçüm Planı}
\label{sec:olcum-plani}

Gelişme raporundaki 20--30 kişilik ölçüm hedefi, sergi takvimi içinde
gerçekleştirilemediği için pratik düzeyde \textbf{küçültülmüş bir test
seti} kullanılmıştır. Test seti üç bileşenden oluşur:

\begin{itemize}[itemsep=2pt]
    \item \textbf{Gallery (kayıtlı):} Proje ekibi (Göktürk Göçen, Ekin
          Ağaoğlu) ve gönüllü olarak katılan toplam 5--8 kişi. Her bir
          birey için kayıt fotoğrafı düzgün aydınlatma ve frontal poz
          ile alınmıştır.
    \item \textbf{Genuine probe seti:} Aynı bireylerin farklı zamanlarda,
          farklı poz/aydınlatmalarda alınmış 3--5 fotoğrafı.
    \item \textbf{Impostor probe seti:} Gallery'de olmayan bireylere ait
          (kamuya açık veri setlerinden seçilmiş) fotoğraflar.
\end{itemize}

Her probe \texttt{far\_frr.py} aracı ile 10 karelik bir burst olarak
sunucuya gönderilmiş; sunucu kalite filtresi, centroid hesabı ve kosinüs
benzerliği üretmiştir. Test üç farklı koşul grubunda tekrarlanmıştır:

\begin{enumerate}[label=(\arabic*),itemsep=2pt]
    \item \textbf{Aydınlatma}: 100~lx (alacakaranlık), 300~lx (ofis), 600~lx (gün ışığı).
    \item \textbf{Mesafe}: 30--60~cm (yakın), 60--90~cm (orta), 90--120~cm (uzak).
    \item \textbf{Pasif canlılık}: gerçek yüz vs.\ ekran tekrar (replay) testleri.
\end{enumerate}

\paragraph{Veri Yönetimi.} Tüm test görüntüleri sergi sonrası
silinmek üzere yalnızca yerel diskte saklanmıştır; gönüllü katılımcılardan
açık rıza alınmıştır. Sunucu tarafında ham görüntü diske yazılmamış,
yalnızca \texttt{far\_frr.py}'nin CSV çıktıları (yüz görüntüsü yok,
similarity skorları var) raporlamaya alınmıştır.

\section{Sonuç Tabloları (Yer Tutucu)}
\label{sec:sonuc-tablolari}

Ölçüm dönemi tamamlandıkça aşağıdaki tablolar gerçek değerlerle
doldurulacaktır. Mevcut sürümde değerler ``--'' ile bırakılmıştır.

\begin{table}[h!]
\centering
\caption{Aydınlatma koşullarına göre FAR/FRR (orta mesafe, $\tau$ = 0.40).}
\label{tab:results-light}
\small
\begin{tabular}{l c c c c}
\toprule
\textbf{Koşul} & \textbf{Genuine probe} & \textbf{Impostor probe} & \textbf{FAR} & \textbf{FRR}\\
\midrule
100~lx (alacakaranlık) & -- & -- & -- & --\\
300~lx (ofis)          & -- & -- & -- & --\\
600~lx (gün ışığı)     & -- & -- & -- & --\\
\bottomrule
\end{tabular}
\end{table}

\begin{table}[h!]
\centering
\caption{Mesafeye göre FAR/FRR (300~lx, $\tau$ = 0.40).}
\label{tab:results-distance}
\small
\begin{tabular}{l c c c c}
\toprule
\textbf{Mesafe}  & \textbf{Genuine probe} & \textbf{Impostor probe} & \textbf{FAR} & \textbf{FRR}\\
\midrule
30--60~cm   & -- & -- & -- & --\\
60--90~cm   & -- & -- & -- & --\\
90--120~cm  & -- & -- & -- & --\\
\bottomrule
\end{tabular}
\end{table}

\begin{table}[h!]
\centering
\caption{Pasif canlılık skoru — gerçek yüz vs.\ fotoğraf replay (300~lx, 60--90~cm).}
\label{tab:results-liveness}
\small
\begin{tabular}{l c c}
\toprule
\textbf{Koşul} & \textbf{Ortalama $\sigma_{emb}$} & \textbf{Standart sapma}\\
\midrule
Canlı yüz (gallery üyeleri) & -- & --\\
Ekran tekrarı (telefonda gösterilen foto) & -- & --\\
\bottomrule
\end{tabular}
\end{table}

\section{Operasyonel Performans}
\label{sec:perf}

Sistemin uçtan uca gecikmesi, \texttt{far\_frr.py} aracının her burst için
ölçtüğü değerlerden hesaplanmıştır. \Cref{tab:results-latency} kategori
bazında ortalama, p50, p95 ve maksimum gecikme değerlerini gösterir.

\begin{table}[h!]
\centering
\caption{Uçtan uca gecikme dağılımı (n adet burst üzerinden, saniye).}
\label{tab:results-latency}
\small
\begin{tabular}{l c c c c c}
\toprule
\textbf{Kategori} & \textbf{n} & \textbf{Ortalama} & \textbf{p50} & \textbf{p95} & \textbf{Maks}\\
\midrule
Tüm burst'ler                          & -- & -- & -- & -- & --\\
Genuine eşleşme                         & -- & -- & -- & -- & --\\
Impostor reddi                          & -- & -- & -- & -- & --\\
Kalite yetersiz (insufficient frames)   & -- & -- & -- & -- & --\\
\bottomrule
\end{tabular}
\end{table}

Gecikme bileşenlerinin yaklaşık dağılımı (varsayılan koşulda, hedef):

\begin{itemize}[itemsep=1pt]
    \item ESP32-CAM görüntü yakalama (10 kare $\times$ 200~ms): $\sim$2.0~s
    \item HTTP POST (Wi-Fi $\geq$~10\,Mbps, 10 kare): $\sim$1.5~s
    \item Sunucu InsightFace çıkarımı (10 kare, CPU): $\sim$2.0~s
    \item Centroid + DB araması: $<$~50~ms
    \item HTTP cevap dönüşü + UART iletim: $<$~200~ms
    \item \textbf{Beklenen toplam:} 5--7~s
\end{itemize}

\section{Eşik Optimizasyonu}

\Cref{ch:gerceklestirme}'de tanımlanan \texttt{far\_frr.py} aracı,
$\tau \in [0, 1]$ üzerinde tam süpürme yapar. EER ve $\mathrm{FAR} \leq 0.01$
kısıtı altındaki operating point, ölçüm tamamlandıktan sonra
aşağıdaki gibi raporlanacaktır:

\begin{itemize}[itemsep=1pt]
    \item \textbf{EER eşiği:} $\tau^\star_{EER}$ = -- (FAR = FRR = --)
    \item \textbf{Operatif eşik} ($\mathrm{FAR} \leq 0.01$): $\tau^\star_{op}$ = -- (FRR = --)
\end{itemize}

Sunucu \texttt{MATCH\_THRESHOLD} ortam değişkeni, ölçüm sonrası optimize
edilen değere göre güncellenip Docker konteyneri yeniden başlatılır.

\section{Donanım Entegrasyon Doğrulaması}

\paragraph{ESP32-CAM $\rightarrow$ EC2.} Uçtan uca burst akışı tezgahta
doğrulanmıştır: ESP32-CAM Wi-Fi STA modunda hotspot'a bağlanmakta, 10
kareyi sunucuya POST etmekte ve JSON yanıtı parse etmektedir. Server
tarafındaki yapılandırılmış JSON günlüğü, gelen burst sayısını,
kalite-geçen kare sayısını ve gecikmeyi izlemekte kullanılmıştır.

\paragraph{STM32 saat ve UART köprüsü.} HSI$\times$PLL ile 216~MHz HCLK
üretilmiş, USART3 VCP'den \texttt{printf} debug çıktısı ve USART6
üzerinden ESP32-CAM'e \texttt{CAPTURE/SCANNING/MATCH/\ldots} satırları
doğrulanmıştır. State machine modülü host-side bir test harness'ında
12 senaryo için ayrıca deterministik olarak doğrulanmıştır
(\Cref{tab:host-tests}).

\paragraph{ESP32-CAM BLE çevresel.} ESP32-CAM cihazı, BLE peripheral
olarak \texttt{TaxiGuard} adıyla yayın yapmakta; \texttt{FFE0} servis
ve \texttt{FFE1} karakteristik üzerinden \emph{notify} akışı
çalışmaktadır. MTU 247 yapılandırması, \texttt{MATCH:<name>;<sim>}
gibi 20 baytı aşan satırların tek pakette iletilmesini garanti eder.

\paragraph{iPhone (TaksiGuvenlik).} SwiftUI uygulaması Login/Register
akışından geçirilmiş, EC2 \texttt{/auth/...} endpoint'leri ile
uçtan uca doğrulanmıştır. Kayıt sonrası otomatik olarak
\texttt{TaxiGuard} BLE çevresel cihazına bağlanmakta, MATCH satırı
geldiğinde \texttt{tel://}<numara> URL'i ile iOS dialer ekranı
otomatik olarak açılmaktadır.

\section{Tartışma}

Sistem mimarisinde alınan üç temel kararın geriye dönüp bakıldığında
hangi pratik faydaları sağladığı şu şekilde değerlendirilebilir:

\begin{itemize}[itemsep=3pt]
    \item \textbf{Kendi sunucumuz vs.\ kapalı kutu API:} İlk bakışta
          AWS Rekognition gibi kapalı kutu bir servis hızlı bir geliştirme
          vaat ediyor olsa da, eşik değerinin kontrolü, model değişiklik
          riski ve KVKK belirsizliği nedeniyle reddedilmiştir. Sergi
          öncesi optimizasyon esnekliği, açık kaynak InsightFace üzerinde
          \texttt{MATCH\_THRESHOLD} ve \texttt{MIN\_QUALITY\_FRAMES}'i
          serbestçe ayarlayabilme imkânını sağlamıştır.
    \item \textbf{Multi-frame vs.\ tek snapshot:} 10 karelik burst, hem
          robust eşleştirme hem ek bir model gerektirmeyen pasif canlılık
          sinyali kazandırmıştır. Maliyeti yalnızca $\sim$5~s'lik
          gecikmedir; bu süre, yolcu binişi başına bir kez gerçekleşen
          bir işlem için makuldür.
    \item \textbf{BLE + telefon vs.\ GSM modülü:} GSM modülü tabanlı
          (SIM800/SIM7600) çözüm, hem donanım maliyeti (en az 250 TL
          ek), hem SIM kart yönetimi hem de operatör servis bedeli gibi
          ekstra yükler getirirdi. Telefonun BLE ile köprüleştirilmesi,
          aynı işlevi sıfır ek altyapıyla sağlamıştır.
    \item \textbf{Entegre BLE vs.\ ayrı modül (HM-10):} İlk Plan B'de yer
          alan HM-10 modülü, ESP32-CAM'in zaten desteklediği BLE
          çevresel arayüzü tarafından yerinden edildi. Bu karar sergi
          öncesinde alındı; tek board üzerinde kamera + Wi-Fi + BLE
          birleşince hem maliyet (~80 TL) hem de ek bir UART hattı
          (STM32 USART2) düştü.
    \item \textbf{iPhone \texttt{tel://} vs.\ Android otomatik çağrı:}
          iOS sandbox tam-otomatik aramaya izin vermez; sürücünün
          dialer ekranında yeşil tuşa tek dokunması istenir
          \cite{apple_tel_url}. Bu, sergi içinde elde olan iPhone~16'nın
          kullanılmasını sağladı ve aynı zamanda yanlış pozitif arama
          koruması olarak savunulabilir bir tasarım sundu.
\end{itemize}

Kısıtlar ve limitler de açıkça raporlanmalıdır:

\begin{itemize}[itemsep=2pt]
    \item Ölçüm seti 5--8 kişilik bir gruba indirgendiği için FAR/FRR
          değerleri istatistiksel olarak sınırlı kapsamlıdır. Sunulan
          metrikler, sistemin \emph{ideal koşullarda yapılabilir} olduğunu
          göstermektedir; geniş ölçekli (n=100+) konuşlandırma için ek
          ölçüm gerekecektir.
    \item Mevcut sunucu HTTP üzerinden çalışmaktadır (Faz 1). Üretim
          ortamı için TLS + token tabanlı kimlik doğrulama zorunludur;
          tasarım buna hazırdır ancak bu tezin değerlendirme kapsamı
          dışındadır.
    \item Yolcunun başörtüsü, gözlük, maske gibi yüz kapayıcı aksesuarlar
          tanımayı etkileyebilmektedir. Bu durum sistem hatası değil,
          ArcFace tabanlı yüz tanımanın doğal bir kısıtıdır ve operasyonel
          süreçte (örnek: birden çok deneme + canlı sürücü onayı) çözülür.
\end{itemize}

\chapter{Sonuç ve Öneriler}
\label{ch:sonuc-oneriler}

\section{Sonuç}

Bu tez kapsamında; taksi gibi araç içi ortamlarda sürücü güvenliğini
artırmaya yönelik, uçtan uca çalışan bir gömülü sistem prototipi
tasarlanmış ve gerçekleştirilmiştir. Sistem; bulut tabanlı yüz tanıma
servisi, kullanıcı kaydı / oturum API'si, ESP32-CAM tabanlı kamera +
Wi-Fi + BLE köprüsü, STM32 tabanlı olay yöneticisi ve iPhone üzerinde
SwiftUI ile yazılmış sürücü uygulamasından oluşmaktadır. Tasarım, üç
temel ilkeye sadık kalmıştır:

\begin{enumerate}[label=(\arabic*),itemsep=2pt]
    \item Sürücü etkileşiminin tek bir butona indirgenmesi,
    \item Görsel/işitsel geri bildirimle dikkat dağıtmayan bir kullanıcı
          deneyimi,
    \item Donanım eki maliyetinin 500~TL altında tutulması.
\end{enumerate}

Bulut tarafında açık kaynak InsightFace kütüphanesinin (buffalo\_l, ArcFace
R100) deploy edilmesi, modelin ve eşik değerinin ekip kontrolünde
kalmasını; bu da KVKK uyumlu bir tasarımı mümkün kılmıştır. Çoklu kare
burst yaklaşımı (10 kare $\times$ 5 FPS), hem doğruluk hem pasif canlılık
açısından tek-kare yaklaşıma kıyasla anlamlı kazanç sağlamıştır. STM32
tarafında HAL'dan soyutlanmış bir durum makinesi modülü yazılması, host
PC üzerinde 12 senaryolu birim testin geçirilmesini ve donanım üzerinde
flash--flaş döngülerinin minimuma indirilmesini sağlamıştır.

Sergi öncesi yapılan tezgah testlerinde:

\begin{itemize}[itemsep=1pt]
    \item EC2 sunucusunun health/check, burst smoke testi ve \texttt{/auth/...}
          mutluluk yolu + hata senaryoları geçmiştir;
    \item ESP32-CAM Wi-Fi STA modunda hotspot'a bağlanıp sunucuya 10
          karelik burst gönderebilmektedir;
    \item ESP32-CAM aynı zamanda BLE çevresel olarak \texttt{TaxiGuard}
          adıyla yayınlanmakta, iPhone uygulaması otomatik olarak bağlanıp
          \texttt{FFE1} karakteristiğinin notify aboneliğini açmaktadır;
    \item STM32 CubeMX projesi açılmış, build geçmiş ve kart üzerinde
          HSI$\times$PLL 216~MHz saat üretimi doğrulanmıştır;
    \item iPhone uygulaması Login/Register akışından geçirilmiş, MATCH
          satırı geldiğinde \texttt{tel://}<numara> dialer'ı otomatik
          olarak açılmaktadır.
\end{itemize}

Sergi haftası içinde tamamlanacak adımlar: küçültülmüş bir test seti
üzerinde FAR/FRR ölçümlerinin alınması, sonuçların tezin
\cref{ch:sonuclar}'üne işlenmesi ve demo öncesi \texttt{emergencyNumber}
değerinin 155'e geri çekilmesi.

\section{Akademik ve Pratik Katkılar}

Çalışmanın özgün katkıları şu şekilde özetlenebilir:

\begin{itemize}[itemsep=2pt]
    \item Açık kaynak yüz tanıma motorunun (InsightFace), küçük ölçekli
          gömülü donanım (ESP32-CAM + STM32) ile entegre edilerek bir
          taksi senaryosuna uyarlanması;
    \item Pahalı GSM modülü yerine telefonun BLE ağ geçidi olarak
          kullanılmasıyla, toplam donanım maliyetinin 500~TL'nin altında
          tutulması;
    \item Çoklu kare embedding standart sapmasının, ek model yükü
          olmadan \emph{pasif canlılık} sinyali olarak kullanılması;
    \item HAL'dan soyutlanmış, host-side birim test edilebilen taşınabilir
          bir STM32 durum makinesi modülü;
    \item Otomatik FAR/FRR ölçüm aracı ile, tezin tablo ve grafiklerinin
          tekrarlanabilir biçimde üretilmesi.
\end{itemize}

\section{Geliştirme Önerileri}

\subsection{Faz 2: HTTPS, ESP-CAM Kimlik Doğrulama ve Token Süresi}

Mevcut sürüm (Faz 1) sunucusu HTTP üzerinden çalışmakta; sergi/test
ortamı için yeterlidir ancak üretim için yetersizdir. Faz 2'ye geçiş için:

\begin{itemize}[itemsep=1pt]
    \item EC2 önüne Caddy veya nginx reverse proxy konulması; Let's Encrypt
          ile otomatik sertifika alınması,
    \item ESP32-CAM \texttt{config.h}'da \texttt{USE\_TLS=1} ve ISRG Root~X1
          sertifika gömülmesi,
    \item Sunucu tarafında \texttt{SHARED\_SECRET} aktif edilerek
          ESP32-CAM ve harness'ın \texttt{X-Shared-Secret} başlığı
          taşıması.
    \item Kullanıcı oturum token'larına süre sınırlaması (örnek 30 gün)
          ve yenileme akışı; mevcut sürümde tokenlar yalnızca çıkış
          işleminde geçersizleştirilmektedir.
    \item iPhone uygulamasında \texttt{NSAllowsArbitraryLoads} bayrağının
          kaldırılması, ATS uyumlu HTTPS'in zorunlu kılınması.
\end{itemize}

Bu adımlar tasarım açısından hazırdır; yaklaşık yarım günlük bir geçiş
işidir.

\subsection{Edge Inference (Kısmi Yerel Tanıma)}

Mevcut tasarım, Wi-Fi gerektirir. Wi-Fi yokken sistem güvenli moda geçer
(\cref{sec:safe-mode}) ancak tanıma çalışmaz. Hibrit bir tasarımda:

\begin{itemize}[itemsep=1pt]
    \item ESP32-S3 veya benzeri bir SoC üzerinde quantized ArcFace-MobileNet
          ile kabaca bir ön-eleme yapılabilir;
    \item Yalnızca yüksek-belirsizlik durumlarında bulut sorgulanır.
\end{itemize}

Bu, hem Wi-Fi yokluğunda kısmi işlevsellik hem de bulut maliyetlerinde
düşüş sağlar; ancak ek bir geliştirme yatırımı gerektirir.

\subsection{Filo Entegrasyonu}

Birden fazla aracın merkezi panelde takibi için MQTT veya WebSocket
tabanlı bir olay panosu eklenebilir. Sunucu zaten yapılandırılmış JSON
günlüğü üretmektedir; bu günlüklerin bir MQTT broker'a yayınlanması
büyük bir mimari değişiklik gerektirmemektedir. Mevcut SQLite kullanıcı
tablosu, her sürücüye atanmış plaka bilgisini zaten tuttuğu için
``hangi araçta hangi şoför var'' raporu küçük bir SQL sorgusu kadar
yakındır.

\subsection{Android İstemcisi}
\label{sec:android-gelecek}

Android tarafında aynı sistem, \texttt{Intent.ACTION\_CALL("tel:155")}
ile tek-dokunuşsuz çalışacak şekilde genişletilebilir
\cite{android2024telephony}. Sürücü tarafında yanlış pozitif arama
sorunu, MATCH onayı için kısa bir geri sayım (örneğin 5~s
``İptal etmek için dokunun'') ile azaltılabilir. Bu, iOS sandbox'ı
gibi dialer ekranı açan bir yaklaşıma kıyasla şoförün ekrana
dokunmadan da emniyet hattı ile bağlantı kurabilmesini sağlar.

\subsection{Aktif Canlılık ve Çoklu Modalite}

Mevcut pasif canlılık (cross-frame embedding $\sigma$) etkili olmakla
birlikte, üst düzey saldırılara karşı (3B maske, yüksek çözünürlüklü
ekran) sınırlı kalabilir. Hareket veya kafa-jest tabanlı aktif challenge
(``başınızı sağa çevirin'') eklenmesi, sürücü deneyimini bozmadan
arka koltukta sessizce gerçekleştirilebilir.

\subsection{Yasal ve Etik Kapsam}

Sistemin sahaya çıkarılması için:

\begin{itemize}[itemsep=1pt]
    \item KVKK kapsamında ``açık rıza'' veya ``kanunlarda öngörülmüş
          hâl'' altında çalışacak hukuki çerçevenin netleştirilmesi,
    \item Aranan kişiler veri tabanına girişin yetkili idari/kolluk
          birimlerince yapılacağına dair operasyonel prosedürün
          tanımlanması,
    \item Yolcunun bilgilendirme yükümlülüğünün (örneğin araç içinde
          görünür yazı veya sesli uyarı) ele alınması gerekmektedir.
\end{itemize}

Bu noktalar mühendislik kapsamı dışındadır ancak sistemin gerçek
konuşlandırmaya yol bulması için ön koşuldur.


% ------------------------------ Kaynakça ---------------------------------
\bibliographystyle{ieeetr}
\bibliography{kaynakca}

% ------------------------------ Ekler ------------------------------------
\appendix
\chapter{Mesaj Protokolleri}
\label{app:protokol}

Bu ek, sistemde tanımlanan tüm UART ve BLE mesaj çerçevelerini tek yerde
listeler. Mesajlar üç hat üzerinde dolaşmaktadır:

\begin{itemize}[itemsep=1pt]
    \item \textbf{UART1} (115200~bps, 8N1): STM32 $\leftrightarrow$ ESP32-CAM
    \item \textbf{UART2} (9600~bps, 8N1): STM32 $\leftrightarrow$ HM-10
    \item \textbf{BLE GATT} (FFE0/FFE1): HM-10 $\leftrightarrow$ Telefon
          (transparent, UART2'nin baytları aynen iletilir)
\end{itemize}

\section{UART1 --- STM32 / ESP32-CAM}

\paragraph{STM32 $\rightarrow$ ESP32-CAM}

\begin{tabular}{l l}
\toprule
\textbf{Komut} & \textbf{Anlam}\\
\midrule
\texttt{CAPTURE\textbackslash n}    & 10-frame burst başlat ve sunucuya gönder\\
\bottomrule
\end{tabular}

\paragraph{ESP32-CAM $\rightarrow$ STM32 (RESULT)}

Burst tamamlandığında veya hata oluştuğunda tek satırlık CSV:

\begin{tabular}{l l}
\toprule
\textbf{Satır}                              & \textbf{Anlam}\\
\midrule
\texttt{1;<name>;<sim>\textbackslash n}     & MATCH; isim ve similarity\\
\texttt{0;;<sim>\textbackslash n}           & NOMATCH; en yakın similarity\\
\texttt{ERR:<tag>\textbackslash n}          & Hata; tag örnek: no\_wifi, http\_500, json\_parse, timeout\\
\bottomrule
\end{tabular}

\section{UART2 ve BLE --- STM32 / HM-10 / Telefon}

HM-10 transparent çalıştığı için STM32'nin UART2'ye yazdığı satırlar BLE
notification olarak telefona iletilir.

\begin{tabular}{l l}
\toprule
\textbf{Satır}                                   & \textbf{Anlam}\\
\midrule
\texttt{SCANNING\textbackslash n}                & TARA başlangıcı\\
\texttt{MATCH:<name>;<sim>\textbackslash n}      & Pozitif eşleşme (telefon 155'i arar)\\
\texttt{NOMATCH\textbackslash n}                 & Negatif burst\\
\texttt{PANIC\textbackslash n}                   & PANİK butonu (telefon 155'i arar)\\
\texttt{NETERR\textbackslash n}                  & Scan timeout veya transport hata\\
\texttt{HB\textbackslash n}                      & 5\,s aralıkla kalp atışı (yalnızca IDLE'da)\\
\bottomrule
\end{tabular}

\section{HTTP --- ESP32-CAM / Sunucu}

\paragraph{İstek.}

\begin{lstlisting}[language=none, basicstyle=\ttfamily\footnotesize]
POST /search HTTP/1.1
Host: 18.192.45.175:8000
Content-Type: application/octet-stream
X-Session-Id: 7c5f9a3b-...-e1
X-Frame-Index: 7
X-Frame-Total: 10
X-Shared-Secret: <token>     # opsiyonel

<JPEG bytes>
\end{lstlisting}

\paragraph{Ara kare yanıtı.}

\begin{lstlisting}[language=none, basicstyle=\ttfamily\footnotesize]
HTTP/1.1 200 OK
Content-Type: application/json

{"status":"continue","frames_received":7,"frames_total":10,"quality_ok_this_frame":true}
\end{lstlisting}

\paragraph{Son kare yanıtı.}

\begin{lstlisting}[language=none, basicstyle=\ttfamily\footnotesize]
HTTP/1.1 200 OK
Content-Type: application/json

{
  "match": true,
  "name": "Ali_Yilmaz",
  "similarity": 0.9437,
  "frames_used": 10,
  "frames_total": 10,
  "liveness_score": 0.014823
}
\end{lstlisting}

\paragraph{Hata yanıtları.}

\begin{tabular}{l l l}
\toprule
\textbf{Kod} & \textbf{Body} & \textbf{Anlam}\\
\midrule
400 & \texttt{\{"error":"bad\_index\_headers"\}}    & X-Frame-* başlıkları geçersiz\\
400 & \texttt{\{"error":"missing\_session\_id"\}}   & X-Session-Id yok veya boş\\
400 & \texttt{\{"error":"bad\_frame\_indices"\}}    & idx > total veya $\leq$ 0\\
400 & \texttt{\{"error":"image\_too\_small"\}}      & JPEG $<$ 1\,KB\\
401 & \texttt{\{"error":"unauthorized"\}}           & X-Shared-Secret eşleşmedi (auth aktif)\\
500 & \texttt{\{"error":"session\_lost"\}}          & Oturum timeout sonrası temizlenmiş\\
\bottomrule
\end{tabular}

\chapter{STM32 Pin Tahsisi ve Bağlantı Şeması}
\label{app:pin-plani}

\section{Pin Tahsis Tablosu (NUCLEO-F767ZI)}

\begin{table}[h!]
\centering
\caption{STM32 NUCLEO-F767ZI çevre birim ve GPIO tahsisleri (Plan B v2).}
\label{tab:pin-tahsisi}
\small
\begin{tabular}{l l l l}
\toprule
\textbf{İşlev} & \textbf{Pin} & \textbf{Bağlantı} & \textbf{Not}\\
\midrule
\multicolumn{4}{l}{\emph{USART6 --- ESP32-CAM iletişimi (115200 bps 8N1)}}\\
TX & PG14 (D1, Arduino başlık)  & ESP32-CAM RX (IO14) & Komut + telefon-yönlü mesaj\\
RX & PG9  (D0, Arduino başlık)  & ESP32-CAM TX (IO13) & RESULT/ERR/HB satırı\\
\midrule
\multicolumn{4}{l}{\emph{USART3 --- ST-LINK VCP debug (115200 bps 8N1)}}\\
TX & PD8 & ST-LINK (Mac terminal)       & \texttt{printf} retarget\\
RX & PD9 & (Kullanılmıyor)               & ---\\
\midrule
\multicolumn{4}{l}{\emph{Buton girişleri (EXTI)}}\\
TARA  & PC13 & B1 USER (kart üstü mavi buton)        & EXTI13, pull-up\\
PANİK & PA0  & Harici buton (sergi düzeneğinde takılı değil) & EXTI0, pull-up\\
\midrule
\multicolumn{4}{l}{\emph{GPIO çıkışları (push-pull)}}\\
LD1 yeşil   & PB0  & NUCLEO onboard            & IDLE / NOMATCH göstergesi\\
LD2 mavi    & PB7  & NUCLEO onboard            & Heartbeat (\~1 Hz) + NETERR blink\\
LD3 kırmızı & PB14 & NUCLEO onboard            & MATCH / PANIC\\
Buzzer      & PD14 & Aktif buzzer 5\,V (transistor sürücülü) & MATCH / PANIC\\
\midrule
\multicolumn{4}{l}{\emph{Kapalı / kullanılmayan çevre birimleri}}\\
ETH      & --- & (Disable)        & DCMI pin çakışması; Plan~B'de DCMI N/A\\
DCMI     & --- & (Disable)        & Kamera ESP32-CAM tarafında\\
USART1   & PA9/PA10 & (Boş)       & Morpho başlığı, kullanılmıyor\\
USART2   & PD5/PD6  & (Boş)       & HM-10 atıldığı için kullanılmıyor\\
\bottomrule
\end{tabular}
\end{table}

\section{Saat Üretimi}

Kart üzerindeki HSE 8~MHz osilatörü tercih edilmemiş; CubeMX'in ürettiği
\texttt{stm32f7xx\_hal\_conf.h} dosyasında \texttt{HSE\_VALUE} makrosu
25~MHz olarak tanımlandığı için (NUCLEO-F767ZI için yanlış varsayılan;
\cite{stm32f767_rm}) HSE üzerinden PLL yapıldığında HAL'in UART BRR
hesabı sapıyor, baud hızı yaklaşık 1/3 oranında bozuluyordu. Onun yerine dahili HSI 16~MHz
osilatörü PLL girişine bağlanır; PLL\_M=8, PLL\_N=216, PLL\_P=2,
PLL\_Q=9 yapılandırmasıyla SYSCLK $=$ HCLK $=$ 216~MHz elde edilir.
APB1 = 54~MHz, APB2 = 108~MHz, USART zaman tabanı APB2'den alınır ve
115200~bps baud hatasız üretilir.

\section{Güç ve Topraklama}

\begin{itemize}
    \item STM32 NUCLEO kartı USB üzerinden (ST-Link konektörü) beslenir.
    \item ESP32-CAM \textbf{ayrı bir USB kaynağından} (kendi dock'undan
          veya bağımsız bir USB şarjdan) beslenir. STM32 kartının 5~V
          pininden besleme \emph{brown-out reset}'ine yol açar; kamera +
          Wi-Fi + BLE eşzamanlı çalışırken pik akım yaklaşık 500~mA'dır.
    \item Tek bir kablo bağlantısı yeterli midir? Hayır: STM32 ile ESP-CAM
          arasında en az üç hat olmalıdır --- USART6 TX, USART6 RX ve
          ortak GND. GND referansı paylaşılmadan UART sinyalleri rastgele
          görünür.
    \item ESP32-CAM tarafında flash sırasında IO0 GND'ye çekilir
          (dock üstündeki düğme aynı işi yapar); sonrasında serbest
          bırakılır.
\end{itemize}

\section{Bring-up Sırası}

Bölüm~\ref{tab:bringup} sergiye gitmeden önce takip edilecek bring-up sırasını
özetler. Her satırın ``başarı kriteri'' sütunu sağlandıktan sonra bir
sonraki satıra geçilir.

\begin{table}[h!]
\centering
\caption{Sergi öncesi donanım bring-up sırası ve başarı kriterleri.}
\label{tab:bringup}
\small
\begin{tabularx}{\linewidth}{l X X}
\toprule
\textbf{Adım} & \textbf{İşlem} & \textbf{Başarı Kriteri}\\
\midrule
1 & STM32'yi CubeIDE'den flash et, USART3 VCP'yi Mac terminal'inde aç (115200) & Boot mesajı + LD2 1~Hz heartbeat + B1 USER basınca SCANNING satırı\\
2 & STM32 USART6 TX'i (D1, PG14) bir USB-TTL ile dinle (115200) & B1 USER basıldığında \texttt{CAPTURE\textbackslash n} ve \texttt{SCANNING\textbackslash n} satırları görünür\\
3 & ESP32-CAM'i dock üzerinde flash et (PlatformIO), kendi dock USB'sinden besle & Seri çıktıda Wi-Fi STA IP'si + ``BLE advertising as TaxiGuard'' mesajı\\
4 & ESP32-CAM'i \texttt{eval/far\_frr.py} ile bağımsız test et & EC2 \texttt{/health} 200, sahte burst için MATCH/NOMATCH cevabı\\
5 & STM32 \texttt{D0/D1/GND} pinlerini ESP32-CAM \texttt{IO13/IO14/GND} pinlerine bağla, iki kart farklı USB güçlerinden çalışsın & B1 USER $\rightarrow$ STM \texttt{CAPTURE} $\rightarrow$ ESP burst $\rightarrow$ STM \texttt{RESULT:...} parse\\
6 & iPhone'da TaksiGuvenlik uygulamasını Xcode ile build et, çalıştır & İlk açılışta Login ekranı; kayıt sonrası \texttt{TaxiGuard}'a otomatik bağlantı\\
7 & B1 USER bas, iPhone'da Scan sekmesini izle & Sırasıyla SCANNING $\rightarrow$ MATCH (veya NOMATCH) durum kartı, MATCH'te \texttt{tel://}<numara> dialer'ı açılır\\
8 & Test seti ile FAR/FRR ölçümü yap (Bölüm~\ref{ch:sonuclar}) & ROC eğrisi + EER raporu üretildi\\
\bottomrule
\end{tabularx}
\end{table}


\end{document}
